% Original analytic teaching diagram, 2026-11-10.
\definecolor{VB}{HTML}{7188A3}\definecolor{VBF}{HTML}{E5EAF0}
\definecolor{VG}{HTML}{829C8B}\definecolor{VGF}{HTML}{E8EEE9}
\definecolor{VY}{HTML}{B4A16F}\definecolor{VYF}{HTML}{F1ECDD}
\definecolor{VR}{HTML}{A7787D}\definecolor{VRF}{HTML}{F0E5E6}
\definecolor{VP}{HTML}{9588A5}\definecolor{VPF}{HTML}{ECE8F0}
\definecolor{VI}{HTML}{3B4252}\definecolor{VM}{HTML}{6B7280}\definecolor{VL}{HTML}{4C4D4F}
\begin{tikzpicture}[x=1mm,y=1mm,>=Stealth,font=\figurefont\fontsize{8}{10}\selectfont,
lab/.style={text=FigureMuted,font=\figurefont\fontsize{7.5}{9}\selectfont},
reg/.style={draw=VB,rounded corners=1.8pt,line width=.8pt},
tt/.style={text=VI,anchor=west},ex/.style={text=VM,anchor=west,font=\figurefont\fontsize{7.5}{9}\selectfont}]
\draw[reg] (0,0) rectangle (108,44);
\draw[reg,fill=VBF!45] (4,4) rectangle (104,37);
\draw[reg,fill=VBF!70] (8,8) rectangle (100,30);
\draw[reg,fill=VBF] (12,12) rectangle (96,23);
\node[tt] at (2,40.5) {连续曲线};
\node[ex] at (50,40.5) {Cantor 函数};
\node[tt] at (6,33.5) {绝对连续曲线};
\node[ex] at (50,33.5) {$2\sqrt{t}$（$t=0$ 处）};
\node[tt] at (10,26.5) {分段连续可微曲线};
\node[ex] at (50,26.5) {切换输入的解};
\node[tt] at (14,17.5) {连续可微曲线};
\node[ex] at (50,17.5) {连续右端的经典解};
\end{tikzpicture}
